\begin{abstract}
The contemporary transition from sidecar-based service meshes to sidecarless, kernel-accelerated designs promises lower per-request overhead, improved node density, and simplified data-plane management. Yet this efficiency shift introduces an unresolved systems tension: the same architectural centralization that reduces the sidecar tax may also increase the risk of cross-tenant performance interference when packet processing and policy enforcement converge in shared node-level components. This proposal investigates that tension through the lens of tail-latency isolation in multi-tenant environments, with particular emphasis on p99 behavior under adversarial and bursty workload conditions. The study focuses on eBPF/XDP-enabled service networking and sidecarless control patterns, including Cilium and Istio Ambient Mesh, to evaluate whether shared kernel-level data planes can deliver both operational efficiency and private performance guarantees. The intended contribution is an empirically grounded methodology and evidence base for assessing isolation quality, noisy-neighbor susceptibility, and mitigation strategies in next-generation cloud-native networking stacks.
\end{abstract}
